\section{Removing completion on the ordinary stratum}\label{h3:sec:ordinary}

The ordinary stratum has two enlargements to remove: adjoining all
Frobenius roots, and allowing unbounded negative powers of \(x\).
Equivariant root retractions give the first comparison.  For the second,
we first prove finiteness for every compact power-series lattice using a
Laurent-tail estimate and a compact stable image.  The intermediate
finiteness argument below uses those lattices and the height-two map to prove the separate bounded-pole
finiteness needed to remove unbounded tails.

\subsection{Coefficient spaces and their topology}

Put $G=P_3$, $A=k[[x,y]]$, and $B_0=A[x^{-1}]$.
For $a,b>0$ write
\[
 v_{a,b}\left(\sum_{i\in\Z,\,j\geq0}c_{ij}x^iy^j\right)
   =\inf_{c_{ij}\ne0}(ai+bj),\qquad v_{a,b}(0)=+\infty.
\]
The ring $B_{\hthreehol}$ consists of the integral-exponent series for which,
for every $a,b>0$ and every real $N$, only finitely many nonzero terms
have $ai+bj\leq N$.  This is precisely convergence on
$0<|x|<1$, $|y|<1$ over trivially valued $k$.
Rational positive $a,b$ give a countable defining family of norms
$\exp(-v_{a,b})$.  Coefficientwise limits of Cauchy sequences show
completeness; Laurent polynomials give a countable dense subset.
Thus $B_{\hthreehol}$ is a Polish additive group.
The algebra $B_{\hthreepk}$ introduced before
\Cref{h3:thm:completed-cohomology} is, equivalently, the completion for
these norms of the union of the finite root levels
$B_{\hthreehol}^{1/p^n}$.  This identification is the comparison of function
algebras in
\Cref{h3:thm:coordinate-comparison,h3:thm:completed-cohomology}.
The countable union of the root Laurent polynomials is dense, so this
complete space is Polish as well.  Frobenius and its inverse are
continuous on $B_{\hthreepk}$, with valuations multiplied and divided
by $p$, respectively.

In particular, \Cref{h3:thm:completed-cohomology} concerns the ordinary
continuous $G$-cochain complex of this space, including continuous
profinite parameters.  We use both its natural constants comparison and its
finiteness assertion for coefficient lines.

The lines needed below are tensor products of integral powers of the
periodicity line and finite determinant-character lines; their duals
are included.  The canonical line is among them by the
canonical-line calculation used in \Cref{h3:lem:invariant-differential}.
Write $\mathcal L$ for such a line, with its original free
$A$-lattice $\mathcal L_A$, and set
\[
 A_d(\mathcal L)=x^{-d}\mathcal L_A\quad(d\in\Z),
 \qquad B_{\hthreehol}(\mathcal L)=B_{\hthreehol}\otimes_A\mathcal L_A,
 \qquad B_{\hthreepk}(\mathcal L)=B_{\hthreepk}\otimes_A\mathcal L_A.
\]
Each $A_d(\mathcal L)$ is $G$-stable and compact, with its
$(x,y)$-adic topology.  This is also its subspace topology in
$B_{\hthreehol}(\mathcal L)$: on a fixed pole lattice, every $v_{a,b}$
is bounded below by a positive multiple of total-degree order minus a
constant, while $v_{1,1}$ recovers that order up to the fixed lattice
shift.  The group preserves $(x,y)$ and carries $x$ to $x$ times a
unit; its action on the chosen line lattice is by units.

For clarity, the notation for algebraic localization always means
\begin{equation}\label{h3:ord:cochain-colimit}
 R\Gamma_{\hthreects}(G,B_0(\mathcal L))
   :=\colim_d C^\bullet_{\hthreects}(G,A_d(\mathcal L)).
\end{equation}
An additional profinite parameter space $T$ means applying
$C(T,-)$ at each displayed finite stage before taking the colimit.
No topology on an unrestricted union of cochain values is being used.

\subsection{Natural root projections}

\begin{lemma}[Equivariant root retractions]\label{h3:ord:root-retractions}
There is a continuous $G$-equivariant $B_0$-linear retraction
$B_0^{1/p}\to B_0$.  Its iterates give compatible retractions
from every finite root level and extend to a continuous retraction
\[
 \rho:B_{\hthreepk}\longrightarrow B_{\hthreehol}
\]
of the natural inclusion.  For every line $\mathcal L$ above, the same
construction gives a continuous equivariant retraction from
$B_{\hthreepk}(\mathcal L)$ to $B_{\hthreehol}(\mathcal L)$.
There are also an integer $e>0$, $q=p^e$, and continuous
$G$-equivariant additive
operators $\mathrm{Fr}$ and $\mathcal P$ on its holomorphic
coefficients such that
\begin{equation}\label{h3:ord:retraction-identity}
 \mathcal P\mathrm{Fr}=1.
\end{equation}
Here $\mathrm{Fr}$ is the $q$-power map in a horizontal frame.
It is invertible on the perfected coefficients, preserves some
$A_{-r}(\mathcal L)$, and $\mathcal P$ preserves $A_D(\mathcal L)$
for every sufficiently large $D$.
All these maps are compatible with constant-field descent.
\end{lemma}

\begin{proof}
By \Cref{h3:lem:ordinary-torsor}, the finite scheme of splittings of the
ordinary connected--\'{e}tale sequence at level $p^\ell$ is a torsor
under
\[
 \underline{\Hom}(\mathcal E[p^\ell],\mathcal C[p^\ell]),
\]
a form of $\mu_{p^\ell}^2$.  Its algebra is
$B_0^{1/p^\ell}$ over $B_0$.
This includes the divisor $y=0$.  A unit $p$-basis on this open is
$x,1+y$; adjoining a root of $1+y$ also adjoins the corresponding
root of $y$.  Thus multiplicative torsor coordinates can be units
throughout the open.  The additive monomial basis $x^iy^j$ used below
only expresses the matrix of the resulting operator; its equivariance
comes from the splitting torsor.
We recall why the construction supplies a retraction without a choice
of a splitting.  After an \'{e}tale cover the group is diagonalizable,
and a comodule is graded by its character group.  Projection onto the
degree-zero summand is linear over the invariant algebra.  Every
comodule map preserves this summand.  Changing the \'{e}tale frame
permutes characters and fixes zero, so the projections descend by
faithfully flat descent.  The invariant subalgebra of a torsor is the
base algebra, and its inclusion is fixed by this projection.
This proves the required retraction.  It is an additive, base-linear
projection; no multiplicativity or preservation of the ideal $(y)$ is
used.  The construction commutes with automorphisms of the ordinary sequence,
hence with $G$ and with coefficient-field automorphisms.  It uses no division by the order of
$\mu_{p^\ell}^2$.

For the untwisted coefficients choose $q=p^e$ with $q$ acting
trivially on $k$, and let $\rho_1:B_0^{1/q}\to B_0$ be the
corresponding retraction.  Set
$\mathcal P(f)=\rho_1(f^{1/q})$ and $\mathrm{Fr}(f)=f^q$.
Then $\mathcal P(g^qf)=g\mathcal P(f)$ and
\eqref{h3:ord:retraction-identity} holds.  On the level
$B_0^{1/q^n}$ define
\begin{equation}\label{h3:ord:finite-root-projection}
 \rho_n(f^{1/q^n})=\mathcal P^n f.
\end{equation}
These maps are $B_0$-linear.  They agree on overlapping levels because
$\mathcal P^{n+1}(f^q)=\mathcal P^nf$.
Thus a single natural root retraction already supplies the whole
compatible system.

Here is also the construction for the actual line twists.  If
$\omega$ denotes the periodicity line, the invariant height-one
section $v_1$ trivializes $\omega^{p-1}$ on this open.
In a free frame it is $x$ times that frame to the power $p-1$,
with the exponent reversed if the dual convention for $\omega$ is
used.  Adjoin a $(p-1)$st root of this section.  This is a finite
\'{e}tale frame cover on the open, and the resulting horizontal frame
has finite tame transition characters.  Include the finite
determinant character and choose $q$ to fix all these characters and
the constants.  The $q$-power map on coefficients in that frame then
descends to the line, as does the base-changed root retraction.
Projection onto the required tame eigensummand is exact, since its
order is prime to $p$.
On the $n$th root level of a line, the intrinsic version of
\eqref{h3:ord:finite-root-projection} is
$\rho_n=\mathcal P^n\mathrm{Fr}^n$; thus its notation does not
require taking a root of a chosen line generator.

In the original rational frame the coefficients in a horizontal frame
have $x$-exponents in a fixed translate $\lambda+\Z$, with
$(q-1)\lambda\in\Z$.  More explicitly, write the horizontal
frame as $h=x^{-\lambda}e_0$, where $e_0$ is the original line
frame, and put $\kappa=(q-1)\lambda$.  Then
\[
 \mathrm{Fr}(fe_0)=f^q x^\kappa e_0.
\]
Thus Frobenius multiplies normalized order by $q$.
For line coefficients we use these normalized Gauss
valuations: relative to the original frame they add the fixed number
$a\lambda$ to $v_{a,b}$.  They define the same topology, and satisfy
$v_{a,b}(\mathrm{Fr}f)=qv_{a,b}(f)$ exactly.
Choosing $r+\lambda\geq0$ makes
$x^r\mathcal L_A$ Frobenius-stable.  This also shows explicitly that
passing to the tame frame introduces only a fixed shift of pole
orders.  The continuity, extension to completion, and assertion about
$A_D$ follow from the estimates proved next; these estimates apply to
the descended eigensummand as well.
\end{proof}

\begin{lemma}[Uniform loss bounds]\label{h3:ord:loss-bounds}
Fix the operators of \Cref{h3:ord:root-retractions} and a free rational
frame of $\mathcal L$, using the normalized Gauss valuations just
specified.  There are constants $c\geq0$ and
$c_{a,b}\geq0$ such that
\begin{align}
 \mathcal P(A_L(\mathcal L))
   &\subseteq A_{L/q+c}(\mathcal L),\label{h3:ord:one-pole-bound}\\
 v_{a,b}(\mathcal Pf)
   &\geq q^{-1}v_{a,b}(f)-c_{a,b}.
   \label{h3:ord:one-norm-bound}
\end{align}
For real indices the pole inequality defines the lattice, so
$A_t=A_{\lfloor t\rfloor}$ in the original integral frame.  With
$C=c/(1-q^{-1})$ and $C_{a,b}=c_{a,b}/(1-q^{-1})$, iteration gives
\begin{align}
 \mathcal P^n(A_L(\mathcal L))
   &\subseteq A_{q^{-n}L+C}(\mathcal L),\label{h3:ord:iterated-poles}\\
 v_{a,b}(\mathcal P^nf)
   &\geq q^{-n}v_{a,b}(f)-C_{a,b}.
   \label{h3:ord:iterated-norms}
\end{align}
The losses for $\rho_n$ in \eqref{h3:ord:finite-root-projection} are
therefore independent of the root level.
\end{lemma}

\begin{proof}
Let $e_0$ be the original free line frame, and temporarily write
$w_{a,b}$ for the unshifted scalar Gauss valuation.  Every scalar
coefficient has the unique decomposition
\[
 f=\sum_{0\leq i,j<q}x^iy^j f_{ij}^{\,q}.
\]
The inverse-semilinear operator is consequently given by a finite
matrix, here a single row:
\[
 \mathcal P(fe_0)=\left(\sum_{0\leq i,j<q} b_{ij}f_{ij}\right)e_0,
 \qquad b_{ij}\in B_0.
\]
The terms of the decomposition have disjoint exponent residue classes,
so
\[
 w_{a,b}(f)=\min_{i,j}\{ai+bj+qw_{a,b}(f_{ij})\}.
\]
Put
\[
 c^0_{a,b}=\max\left(0,
      \max_{i,j}\left\{\frac{ai+bj}{q}-w_{a,b}(b_{ij})\right\}\right).
\]
The scalar output has valuation at least
$q^{-1}w_{a,b}(f)-c^0_{a,b}$.
Since $v_{a,b}(fe_0)=w_{a,b}(f)+a\lambda$, its normalized
valuation is at least
\[
 q^{-1}v_{a,b}(fe_0)-c^0_{a,b}+a\lambda(1-q^{-1}).
\]
Thus \eqref{h3:ord:one-norm-bound} holds, for example, with
$c_{a,b}=c^0_{a,b}+a|\lambda|(1-q^{-1})$.
All $b_{ij}$ have one common finite pole bound.
For $f\in A_L$, its components have pole at most
$L/q+(q-1)/q$.  This proves \eqref{h3:ord:one-pole-bound} with a
constant independent of $L$.
Each iteration divides the previous loss by $q$, giving the displayed
geometric sums and proving \eqref{h3:ord:iterated-poles} and
\eqref{h3:ord:iterated-norms}.

The component decomposition and finite matrix are continuous on
$B_{\hthreehol}$: taking an exponent residue class and then a root
preserves convergence in every defining norm.  On a root level the
last inequality reads
\[
 v_{a,b}(\rho_n h)\geq v_{a,b}(h)-C_{a,b}.
\]
Thus the compatible maps extend continuously from the dense root
union to $B_{\hthreepk}$, with values in the complete space $B_{\hthreehol}$.
Their common restriction is the identity on $B_{\hthreehol}$.
Finally $D>C$ implies $D/q+c<D$, proving stability of $A_D$.
The proof on the tame frame cover uses its finitely many character
components and the same finite matrix calculation.  Descent therefore
retains both estimates for every specified twist.
\end{proof}

\begin{proposition}[Holomorphic comparison]\label{h3:prop:holomorphic-comparison}
For every specified coefficient line the natural map
\[
 R\Gamma_{\hthreects}(G,B_{\hthreehol}(\mathcal L))
   \longrightarrow R\Gamma_{\hthreects}(G,B_{\hthreepk}(\mathcal L))
\]
is a quasi-isomorphism.  Its cohomology groups are finite, and the
chosen Frobenius power acts invertibly on them.  For a profinite
parameter space $T$, the comparison and Frobenius invertibility persist,
and the cohomology is naturally
\[
 C\bigl(T,H^a_{\hthreects}(G,B_{\hthreehol}(\mathcal L))\bigr),
\]
where the finite point-value cohomology group has its discrete topology.
Finiteness is asserted for the cohomology without parameters.
The comparison is compatible with coefficient-field descent.
For the untwisted line it identifies
the natural constants comparison of \Cref{h3:prop:ordinary-constants}.
\end{proposition}

\begin{proof}
The retraction induces a split injection on cohomology into the finite
groups of \Cref{h3:thm:completed-cohomology}.  Hence holomorphic
cohomology is finite before any bounded-pole assertion is used.
The identity $\mathcal P\mathrm{Fr}=1$ makes Frobenius injective
on these finite groups, hence bijective.  Both cochain complexes are
Polish, so their finite cohomology groups are discrete by
\Cref{h3:co:strictness}.

For completeness, consider a completed class, rather than merely a
class in the algebraic root union.  On perfected coefficients define
\[
 R_\ell=\mathrm{Fr}^{-\ell}\rho\mathrm{Fr}^{\ell}:
 B_{\hthreepk}(\mathcal L)\longrightarrow
 B_{\hthreehol}(\mathcal L)^{1/q^\ell}.
\]
This is a projection onto the indicated root level and satisfies
\begin{equation}\label{h3:ord:root-approximation-bound}
 v_{a,b}(R_\ell f)\geq v_{a,b}(f)-q^{-\ell}C_{a,b}.
\end{equation}
It converges to the identity uniformly on compact subsets in every
defining norm.  Indeed, approximate a compact subset, at any prescribed
accuracy in finitely many norms, by finitely many elements at a common
finite root level.  For every later $\ell$, $R_\ell$ fixes these
approximants.  The bound \eqref{h3:ord:root-approximation-bound} controls
the errors uniformly; increasing the accuracy proves the assertion.

If $z$ is a continuous completed cocycle, its image is compact.
Thus $R_\ell z\to z$ in the cochain topology.  These are cocycles
because the projections are equivariant.  Discreteness of completed
cohomology gives $[R_\ell z]=[z]$ for all sufficiently large $\ell$.
Every completed class therefore comes from a finite root level.
Under Frobenius transport, the image of this level is the image of
holomorphic cohomology, because Frobenius is already invertible on
that image.  This proves surjectivity and hence the comparison.

The finite-cohomology lifting assertion of \Cref{h3:co:parameters}
now applies to these two Polish complexes.  It proves the comparison
with any profinite parameters.  All the maps used to define the
comparison are the natural inclusions; the retractions prove their
properties.  In particular the comparison preserves products and the
original constants maps.  Naturality of the torsor construction also
proves its coefficient-field equivariance.
\end{proof}

\subsection{Uniform Laurent tails and compact stable images}

Let $T_D$ discard the terms of pole order at most $D$, retaining only
$x^iy^j$ with $i<-D$, in a chosen rational frame.  It is a continuous
linear operator on $B_{\hthreehol}(\mathcal L)$; its complementary
projection has image $A_D(\mathcal L)$.  The frame shifts in
\Cref{h3:ord:root-retractions} can be absorbed in the constants of
\Cref{h3:ord:loss-bounds}.

\begin{lemma}[Tail contraction on compact families]\label{h3:ord:tail-contraction}
Choose an integer $D>C$ in \Cref{h3:ord:loss-bounds}.
For every compact subset $K\subset B_{\hthreehol}(\mathcal L)$,
\[
 T_D\mathcal P^nf\longrightarrow0
 \quad\hbox{uniformly for }f\in K
\]
in every defining Gauss norm.
\end{lemma}

\begin{proof}
Choose $0<\delta<D-C$ and set $L_n=\lfloor\delta q^n\rfloor$.
The pole estimate implies
\[
 \mathcal P^n(1-T_{L_n})f\in A_D(\mathcal L),\qquad
 T_D\mathcal P^nf=T_D\mathcal P^nT_{L_n}f.
\]
Fix $a,b>0$.  For any $R>0$, compactness gives a finite lower bound
\[
 m_R:=\inf_{f\in K}v_{a+R,b}(f)>-\infty.
\]
A monomial with pole $i>L$ has normalized $v_{a,b}$-value equal
to its $v_{a+R,b}$-value plus $R(i-\lambda)$.  Consequently
\[
 \inf_{f\in K}v_{a,b}(T_{L_n}f)\geq m_R+R(L_n-\lambda).
\]
Truncation cannot lower a Gauss valuation.  By the iterated norm bound,
\begin{equation}\label{h3:ord:quantitative-tail}
 \inf_{f\in K}v_{a,b}(T_D\mathcal P^nf)
   \geq q^{-n}(m_R-R\lambda)+Rq^{-n}L_n-C_{a,b}.
\end{equation}
The lower limit is at least $R\delta-C_{a,b}$.  Since $R$ can
be arbitrarily large, it is $+\infty$.  This includes $\lambda=0$
for the untwisted line.
\end{proof}

\begin{lemma}[The compact kernel and its stable image]\label{h3:ord:stable-image}
For $D$ as above and every $a\geq0$, put
\[
 M=H^a_{\hthreects}(G,A_D(\mathcal L)),\qquad
 \phi:M\longrightarrow H^a_{\hthreects}(G,B_{\hthreehol}(\mathcal L)).
\]
Then $M$ is compact Hausdorff, and
$\mathcal P^nz\to0$ for every $z\in\ker\phi$.
Moreover its stable image
\[
 M_\infty=\bigcap_{n\geq0}\mathcal P^nM
\]
is finite.
\end{lemma}

\begin{proof}
Compactness, Hausdorffness and the asserted cohomology topology follow
from the finite completed resolution in \Cref{h3:co:compact-resolution}.
The continuous map $\phi$ has finite discrete target by
\Cref{h3:prop:holomorphic-comparison}, so $K=\ker\phi$ is compact.
For $a=0$ it is zero, since the coefficient inclusion is injective.
For $a>0$, represent $z\in K$ by an $A_D$-valued cocycle $z_0$
and write $z_0=\partial b$ with a holomorphic continuous cochain $b$.
Set
\[
 b_n=(1-T_D)\mathcal P^nb,\qquad
 e_n=\mathcal P^nz_0-\partial b_n
      =\partial(T_D\mathcal P^nb).
\]
The first expression shows that $e_n$ is an $A_D$-valued cocycle.
The image of $b$ is compact, so \Cref{h3:ord:tail-contraction} makes
$T_D\mathcal P^nb$ tend uniformly to zero.
The group differential preserves this convergence in total-degree
order: every group element preserves $(x,y)$, sends $x$ to $x$ times
a unit and sends a line generator to a unit multiple.  Thus its
action does not decrease total-degree order in a fixed rational
frame.  The finite sum defining $\partial$ has the same property.
It follows that $e_n\to0$ in the compact lattice topology, and
$\mathcal P^nz=[e_n]\to0$ in $M$.

We give the uniformity step explicitly.  For an open additive subgroup
$U$ of $K$, let
\[
 K_N(U)=\bigcap_{n\geq N}\{z\in K:\mathcal P^nz\in U\}.
\]
These are increasing closed subgroups covering $K$ by the convergence
just proved.  The Baire theorem makes one of them have interior; as a
subgroup it is open and has finite index.  Choose representatives for
its finitely many cosets.  Each representative has all sufficiently
late iterates in $U$.  Taking the largest of these finitely many
bounds gives $\mathcal P^nK\subseteq U$ for all sufficiently
large $n$.  Hence
\begin{equation}\label{h3:ord:kernel-intersection}
 \bigcap_{n\geq0}\mathcal P^nK=0.
\end{equation}

The operator $\mathcal P$ is bijective on holomorphic cohomology:
it is the inverse of Frobenius there by finiteness and
\eqref{h3:ord:retraction-identity}.
If $z\in M_\infty\cap K$, choose, for each $n$, an $m_n\in M$
with $\mathcal P^nm_n=z$.  Applying $\phi$ and using this
bijectivity gives $m_n\in K$.  Equation
\eqref{h3:ord:kernel-intersection} therefore forces $z=0$.
Thus $M_\infty$ injects into the finite holomorphic cohomology group,
which proves its finiteness.
\end{proof}

\begin{proposition}[Finiteness of all ordinary lattices]
\label{h3:prop:ordinary-lattice-finiteness}
For every $a\geq0$, $d\in\Z$, and coefficient line $\mathcal L$
specified above, the group
\[
 H^a_{\hthreects}(P_3,A_d(\mathcal L))
\]
is finite.  This includes all periodicity, determinant, conormal, and
canonical-line twists and their inverses, with the original
$(x,y)$-adic lattice topology.  The conclusion uses no finiteness
or cohomology comparison for $B_0$.
\end{proposition}

\begin{proof}
Each adjacent quotient
$A_j(\mathcal L)/A_{j-1}(\mathcal L)$ is a free
$k[[y]]$-line.  Its action is the restriction of $\mathcal L$
tensored with the appropriate power of the conormal of $x=0$.
That conormal is a periodicity power.  Thus
\Cref{h3:co:base-lattices} says that every such quotient has finite
$G$-cohomology.  Finite strips of adjacent quotients have finite
cohomology as well, by their finite filtration.  The sequences are
exact on continuous cochains: the finite Laurent-coordinate
truncations give continuous nonequivariant sections.

Choose $A_{-r}(\mathcal L)\subseteq A_D(\mathcal L)$ with the
former Frobenius-stable and the latter $\mathcal P$-stable, as in
\Cref{h3:ord:root-retractions}.  If
$z\in H^a(G,A_{-r}(\mathcal L))$, then its image in $M$ satisfies
\[
 \iota(z)=\mathcal P^n\iota(\mathrm{Fr}^nz)
 \quad\hbox{for every }n.
\]
It therefore belongs to the finite subgroup $M_\infty$ of
\Cref{h3:ord:stable-image}.  The kernel of this map is a quotient of
$H^{a-1}(G,A_D/A_{-r})$, which is finite by the strip calculation
(and is zero for $a=0$).
Hence $H^a(G,A_{-r})$ is finite.  The same strip sequence then
proves finiteness for $A_D$; further finite strips reach every
$A_d$, in either direction.  All twists have been retained throughout.
\end{proof}

\subsection{Intermediate finiteness and bounded poles}
\label{h3:sec:intermediate-finiteness}

The compact lattice calculation now supplies a finite chromatic
abutment.  We use it to bound the algebraic bounded-pole cohomology
before invoking the ordinary comparison.

Write \(M_3\) for the fiber of \(L_3\to L_2\).

\begin{lemma}\label{h3:des:monochromatic-finiteness}
The groups
\[
 \pi_j(T/p),\qquad \pi_j(M_3S/p),\qquad
 \pi_j(X/p),\qquad \pi_j(L_1X/p)
\]
are finite for every integer \(j\).  This assertion uses
\Cref{h3:prop:ordinary-lattice-finiteness} but does not use
\Cref{h3:thm:ordinary-comparison}.
\end{lemma}

\begin{proof}
\emph{The height-three local term \(T/p\).}
Apply the finite Moore quotient to the relative object of
\Cref{h3:prop:exact-descent}.  Its abutment is
\(\pi_*(R_{3,k}/p)\), and its coefficient in internal degree
\(2j\) is \(A\otimes\omega^j\); odd coefficients vanish.
These compact modules have finite \(P_3\)-cohomology by
\Cref{h3:prop:ordinary-lattice-finiteness}.  They are complete and
separated with invariant open neighborhoods, so the reduced
\(\F_p[[P_3]]\)-resolution of \Cref{h3:co:compact-resolution},
which has length nine, computes their continuous cochains.  Hence only
columns \(0\leq s\leq9\) occur.  The bounded convergence from
\Cref{h3:prop:exact-descent} gives finitely many finite associated-graded
groups in each total degree, so \(\pi_*(R_{3,k}/p)\) is degreewise
finite.  The underlying equivalence
\(R_{3,k}/p\simeq(T/p)^{\vee [k:\F_p]}\) then gives finiteness for
\(T/p\) over \(S\).

\emph{The monochromatic term \(M_3S/p\).}
The chromatic fracture square gives natural equivalences
\[
 M_3S\simeq M_3T\simeq\fib(T\longrightarrow L_2T).
\]
The first identity is also \cite[Theorem~6.19]{HS99}.
Put $E=E_3(k)$, $\widetilde A=\pi_0E=W(k)[[x,y]]$ and
$I=(p,x,y)$. The module-spectrum form of chromatic localization is
needed here. For every ordinary $E$-module $N$, including a completed
continuous-function term, there is a natural fiber sequence
\[
 \operatorname{Kos}_E(I)\otimes_E N\longrightarrow N
       \longrightarrow L_2N,
 \qquad
 \operatorname{Kos}_E(I)=
 \bigotimes_{z\in\{p,x,y\}}\fib(E\longrightarrow E[z^{-1}]).
\]
This is \cite[Proposition~3.4(1) and Remark~3.5]{BarthelStapleton2016},
with its support functor described in Lemmas~2.3 and~2.5 of that paper.
The statement is about the underlying chromatic localization of
arbitrary module spectra; it imposes no flatness or finite-generation
condition. Since $E$ is $E(3)$-local and $L_3$ is smashing, every
$E$-module is $E(3)$-local. Thus the first term is $M_3N$.
In particular it applies to the actual completed Amitsur spectrum
$N_Q=\mathcal N_3^q$ for $Q=P_3^q$, equipped with
\eqref{h3:des:relative-cooperations} and viewed as an ordinary
$E$-module; no completed tensor product is replaced by an ordinary one.

The spectrum $N_Q/p$ has zero $p$-inversion, so its $p$-support
factor is the whole spectrum. The remaining two factors give the
finite \v Cech support complex for $(x,y)$. Its homotopy spectral
sequence is
\[
 H^s_{(x,y)}(\pi_t(N_Q/p))
       \Longrightarrow \pi_{t-s}(M_3(N_Q/p)),\qquad 0\leq s\leq2.
\]
By \Cref{h3:des:flat-functions}, the even coefficient modules of
$N_Q/p$ are $C_{\hthreects}(Q,A)\otimes\omega^{t/2}$,
where $A=k[[x,y]]$, and are flat over $A$; the odd ones vanish.
The regular sequence $x,y$ therefore makes the displayed support
cohomology vanish except for $s=2$. For $Q$ a point it is
\[
 H^2_{(x,y)}A
 =A[x^{-1},y^{-1}]/(A[x^{-1}]+A[y^{-1}]).
\]
Thus the shift is exactly two. The finite support complex introduces
no further convergence issue.

These identifications also respect the action-cobar maps.  All faces
other than the first, and all degeneracies, are $E$-linear; the first face uses the
continuous stabilizer action. The invariant ideal $I$ is preserved
even for this varying family: in
$C_{\hthreects}(P_3^{q+1},\widetilde A)$ one has
$(p,g_1(x),g_1(y))=(p,x,y)$. Indeed, the continuous monomial-division
operations used in \Cref{h3:des:flat-functions} express any element
of $I$ as a continuous sum of $p,x,y$ multiples. Apply them to
$g_1^{-1}(x)$ and $g_1^{-1}(y)$, then apply $g_1$; joint continuity
gives continuous coefficients expressing $x,y$ in the image ideal.
The reverse containment follows in the same way. For precision, let $I_q$ be the image of $I$ in $\pi_0N_Q$.
Base change of the displayed finite fiber construction from $E$ to
$N_Q$ identifies the ordinary $E$-support term with
$\operatorname{Kos}_{N_Q}(I_q)$.  Along an actual coface ring map
$N_Q\to N_{Q'}$, its further base change is the construction for the
image generators in $\pi_0N_{Q'}$.  The support functor on ordinary
$N_{Q'}$-modules is the canonical right adjoint for modules supported
at that finitely generated ideal, so it depends on the ideal, not on
the chosen generators.  The equality of ideals just proved identifies
the first-face image with the target support ideal.  Uniqueness of this
right adjoint makes the resulting comparisons compatible with the first face and with compositions.  No individual generator
is assumed fixed by the action. The calculation therefore applies to the
whole completed Amitsur object, and \Cref{h3:prop:exact-descent}
permits its termwise use before totalization.


For a continuous function term the same quotient is
\begin{equation}\label{h3:des:principal-parts-parameters}
 H^2_{(x,y)}C_{\hthreects}(Q,A)
     =C_{\mathrm{lc}}(Q,H^2_{(x,y)}A),
\end{equation}
where the target coefficient module is discrete.
To verify this, express \(H^2_{(x,y)}A\) as the colimit of
\(A/(x^a,y^a)\), with transition multiplication by \(xy\).
At each stage, quotienting continuous functions gives continuous
functions to the finite discrete quotient, using coefficient
truncation to lift them.  A locally constant map from compact
\(Q\) has finite image, so one common stage represents it.
This proves \eqref{h3:des:principal-parts-parameters} with the
required denominator and parameter bounds.

Residue pairing, followed by \(\hthreeTr_{k/\F_p}\), identifies the
continuous dual of
\(H^2_{(x,y)}(A\otimes\omega^j)\) with the compact module
\[
 M=A\otimes\omega^{-j}\otimes
       {\textstyle\bigwedge^2}\Omega^1_{A/k,\hthreects}.
\]
Concretely the pairing is the coefficient of \(x^{-1}y^{-1}\)
in a product with a two-form.  It is perfect on finite
principal-parts quotients and hence on their colimit against
the inverse-limit power-series module.  Its change-of-variables
formula proves equivariance.  The canonical-line identity of
\Cref{h3:lem:invariant-differential} makes this an admissible twist
in \Cref{h3:prop:ordinary-lattice-finiteness}.
The principal-parts module is discrete and \(M\) is compact.
The reversed compact/discrete duality
\eqref{h3:co:compact-duality}, applied to \(P_3\) in dimension nine,
therefore gives
\[
 H^s_{\hthreects}\bigl(P_3,H^2_{(x,y)}(A\otimes\omega^j)\bigr)
   \cong H^{9-s}_{\hthreects}(P_3,M)^\vee.
\]
The right side is finite by
\Cref{h3:prop:ordinary-lattice-finiteness}, so the principal-parts
cohomology is finite in every degree.
The relative descent spectral sequence for
\(M_3(R_{3,k}/p)\) is bounded in columns by the same reduced
length-nine resolution and has finite entries in every total degree.
It therefore has degreewise finite abutment.  Its underlying spectrum
is a sum of \([k:\F_p]\) copies of \(M_3(T/p)\), because the exact
functor \(M_3(-/p)\) preserves the finite free decomposition of
\(R_{3,k}\).  Using \(M_3T\simeq M_3S\) proves finiteness of
\(\pi_*(M_3S/p)\).

\emph{The overlap \(X/p\).}
The fiber sequence \(M_3S/p\to T/p\to X/p\) gives finiteness
for \(X/p\).

\emph{The height-one localization \(L_1X/p\).}
By \Cref{h3:thm:height-two-test},
\(L_{K(2)}X/p\) is the sum of two shifts of \(R_2/p\).
The completion map identifies \(S/p\) with the finite torsion Moore
spectrum \(\mathbb S_{(p)}/p\), and exactness of \(L_{K(2)}\) gives
\[
 R_2/p\simeq L_{K(2)}(S/p).
\]
For \(p\geq5\), \cite[Theorem~15.1]{HS99} applied at height two to this
finite torsion spectrum says that every integer-graded homotopy group is
finite; see also \cite[p.~49(c)]{BB}.

For the other height-two boundary put \(E=L_1R_2\).  The completion
map from the \(p\)-local sphere to \(S\) has rational cofiber, so it is
a \(K(2)\)-equivalence.  Thus \(R_2\) is the \(K(2)\)-local sphere in
the group calculation of \cite[Remark~7.8]{Behrens2012}.  We use the
following features of that formula.  Write \(d_v=|v_1|=2(p-1)\).
Besides finitely many shifted copies of \(\Z_p\), \(\Q_p\), and
\(\Q/\Z_{(p)}\), the formula contains finite cyclic \(p\)-groups with
generators \(1_{sp^n/(n+1)}\), where \(n\geq0\), \(p\nmid s\), and the
internal degree is \(d_vsp^n\).  Its exterior factor on the degree-minus-one
class denoted here by \(\zeta_{\mathrm B}\) adds only two shifts.  This
notation distinguishes that height-two class from the map \(\zeta\)
constructed above.

For a fixed nonzero internal degree, division by \(d_v\), when possible,
determines at most one integer of the form \(sp^n\) with \(p\nmid s\):
\(n\) is its \(p\)-adic valuation and \(s\) its remaining \(p\)-unit
integer.  The cyclic family therefore contributes at most one summand
in each such degree before the two exterior shifts.  Multiplication by
\(p\) has finite kernel and cokernel on each listed group: these pairs
are \((0,\F_p)\), \((0,0)\), and \((\F_p,0)\) on \(\Z_p\), \(\Q_p\),
and \(\Q/\Z_{(p)}\), respectively, and they are finite on a finite
cyclic \(p\)-group.  Hence both \(\pi_j(E)[p]\) and \(\pi_j(E)/p\) are
finite for every integer \(j\).  The Moore cofiber sequence gives
\[
 0\longrightarrow\pi_j(E)/p\longrightarrow\pi_j(E/p)
   \longrightarrow\pi_{j-1}(E)[p]\longrightarrow0.
\]
It follows that \(\pi_j(E/p)\) is finite in every degree.  Exactness
identifies \(E/p\simeq L_1(R_2/p)\), the boundary needed here.
The local height-two equivalence then identifies \(L_1L_{K(2)}(X/p)\)
with two shifts of \(E/p\), so that corner is degreewise finite.
Finally, apply the height-two fracture square to \(X/p\).
Its other three corners are degreewise finite, so its long
exact sequence proves finiteness of \(L_1X/p\).
No bounded-pole comparison has been used in this argument.
\end{proof}

\begin{proposition}\label{h3:prop:bounded-pole-finiteness}
For every \(a\geq0\), the group
\(H^a_{\hthreects}(P_3,B_0)\), with the prescribed bounded-pole
cochain convention, is finite.
\end{proposition}

\begin{proof}
Apply \(L_1(-/p)\) to the completed relative descent object.
We first identify this exact test on an actual term \(N_Q=\mathcal N_3^q\).
For odd \(p\), an Adams self-map of the finite Moore spectrum has degree
\(2(p-1)\) and acts nontrivially in \(K(1)\)-homology
\cite[Example~2.4.1(ii)]{RavenelNilpotence}.  Its action on
\(BP_*(S/p)=BP_*/p\) is multiplication by an element of degree
\(2(p-1)\), and that graded group is \(\F_p v_1\).  This even-degree
coefficient detects the entire \(BP\)-module map: applying maps into
\(BP/p\) to \(BP\xrightarrow p BP\to BP/p\) shows that the only
possible extra term is \(\pi_{2(p-1)+1}(BP/p)=0\).  If the coefficient
were zero, the \(BP\)-module map would be null, and so would its base
change to \(K(1)\), contrary to the chosen Adams map.  Its coefficient
is therefore a unit multiple of \(v_1\).  Rescaling the self-map by an
integer unit makes its \(BP\)-Hurewicz action exactly \(v_1\).
The height-three orientation sends this element to \(v_1=xU^{p-1}\)
modulo \(p\).  This unit rescaling does not change the telescope:
powers of the unit give an isomorphism between the two defining towers.

The height-one telescope theorem for this finite Moore spectrum is the
consequence of \cite[Theorem~4.11 and Corollary~4.12]{Miller1981}
identified as the \(n=1\) case in
\cite[Section~7.5, after Conjecture~7.5.5]{RavenelNilpotence}.
Since \(L_1\) is smashing, it gives, for the already formed ordinary
module \(N_Q\),
\[
 L_1(N_Q/p)\simeq N_Q\wedge L_1(S/p)
       \simeq N_Q\wedge\operatorname{Tel}(v_1:S/p\to\Sigma^{-2(p-1)}S/p).
\]
Thus the test is a termwise ordinary telescope of the same finite Moore
self-map.  Homotopy groups commute with that telescope.  In internal
degree zero the resulting coefficient is exactly
\[
 \colim_d x^{-d}C_{\hthreects}(P_3^q,A).
\]
Its other coefficients are the periodic translates.  The actual action
preserves every bounded-pole lattice modulo \(p\), because
\(g(x)/x\) is a unit in \(A\).  This is the prescribed cochain complex
of \(B_0\); no completion of the algebraic localization is taken.

The central tame units act trivially on \(A\) and by their
scalar powers on \(U\).  Averaging over \(\mu_{p-1}\) kills
internal degrees not divisible by \(2(p-1)\).
Modulo \(p\), \(v_1\) is invariant, so multiplication by
its powers identifies all the remaining coefficient complexes.
For each compact lattice \(x^{-d}A\), the reduced finite
\(\F_p[[P_3]]\)-resolution of \Cref{h3:co:compact-resolution}
computes its continuous cochains and has length nine.  Taking the exact
filtered colimit of these complexes proves the same column bound for
this prescribed bounded-pole union; it does not apply a compact-module
theorem directly to the union.  Thus the spectral sequence has
\[
 E_2^{a,\,2b(p-1)}
       =H^a_{\hthreects}(P_3,B_0)\,v_1^b,\qquad
 0\leq a\leq9,
\]
and no other nonzero internal degrees.  The abutment is the relative
one, \(\pi_*L_1(R_{3,k}/p)\).  If \(m=[k:\F_p]\), the underlying
finite free decomposition gives
\[
 L_1(R_{3,k}/p)\simeq L_1(T/p)^{\vee m}
       \simeq L_1(X/p)^{\vee m}.
\]
Here \(L_1L_2=L_1\), and all functors involved preserve finite sums.
Thus this relative abutment is degreewise finite by
\Cref{h3:des:monochromatic-finiteness}.  The error-tower bound in
\Cref{h3:prop:exact-descent} and the finite column range give the
bounded filtration whose associated graded is \(E_\infty\).

The differential has bidegree \((r,r-1)\), so a nonzero
differential requires \(2(p-1)\mid r-1\).
For \(p\geq7\), no such \(r\geq2\) fits the column range;
the finite abutment then gives the result directly.
For \(p=5\), the only possible differential is
\[
 d_9:E_9^{0,t}\longrightarrow E_9^{9,t+8}.
\]
Its source is finite.  Indeed, coefficient inclusion gives
\[
 H^0(P_3,B_0)\lhook\joinrel\longrightarrow
 H^0(P_3,B_{\hthreehol}),
\]
and the target is finite by
\Cref{h3:prop:holomorphic-comparison}.
All earlier differentials vanish by the displayed sparsity,
and hence the \(d_9\)-image is finite.  Column zero is already finite
by this invariants argument.  Columns one through eight are unchanged
and are finite from their \(E_\infty\) groups.  In column \(9\),
\(E_2^{9,t}=E_9^{9,t}\) is an extension of its finite
\(E_\infty\)-quotient by that finite image, and is finite.
This proves the proposition also at five, without asserting
that \(d_9\) vanishes.
\end{proof}

\subsection{The natural ordinary comparison}
\label{h3:sec:ordinary-comparison}

The preceding proposition supplies the finiteness needed for the second
decompletion.  We now remove unbounded Laurent tails by combining that
finiteness with the contraction already proved on every fixed quotient.

\begin{theorem}[Natural ordinary comparison]\label{h3:thm:ordinary-comparison}
The natural coefficient maps induce quasi-isomorphisms
\begin{equation}\label{h3:ord:natural-comparison}
 \colim_d C^\bullet_{\hthreects}(P_3,x^{-d}A)
 \longrightarrow C^\bullet_{\hthreects}(P_3,B_{\hthreehol})
 \longrightarrow C^\bullet_{\hthreects}(P_3,B_{\hthreepk}).
\end{equation}
They preserve products, constants, and coefficient-field descent, and
remain quasi-isomorphisms with any profinite space of continuous
parameters in the convention \eqref{h3:ord:cochain-colimit}.
Every Frobenius power, including the $p$th-power map, acts invertibly
on these untwisted cohomology groups.
Consequently there is the natural identification
\[
 H^*_{\hthreects}(P_3,B_0)
   \cong H^*_{\hthreects}(P_1\times\GL_2(\Z_p),k).
\]
\end{theorem}

\begin{proof}
By \Cref{h3:prop:bounded-pole-finiteness}, the cohomology of the
left-hand complex in \eqref{h3:ord:natural-comparison} is finite.
The lattice, height-two, and spectral arguments establishing this input
have already been completed.

First keep $D>C$ fixed.  The quotient
$Q_D=B_{\hthreehol}/A_D$ is Polish and has the continuous tail section
$T_D$.  Hence
\[
 0\longrightarrow C^\bullet_{\hthreects}(G,A_D)
 \longrightarrow C^\bullet_{\hthreects}(G,B_{\hthreehol})
 \longrightarrow C^\bullet_{\hthreects}(G,Q_D)\longrightarrow0
\]
is exact degreewise.  Its long exact sequence, together with
\Cref{h3:prop:holomorphic-comparison,h3:prop:ordinary-lattice-finiteness},
shows that $H^a(G,Q_D)$ is finite.  It is discrete by
\Cref{h3:co:strictness}.  Since $\mathcal P$ preserves $A_D$, it
acts on this fixed quotient.
Represent a quotient cocycle by its continuous holomorphic tail
section $f$.  Its iterates are represented by
$T_D\mathcal P^nf$, which converge to zero by
\Cref{h3:ord:tail-contraction}, uniformly on the compact cochain domain.
These are quotient cocycles.  Discreteness therefore implies that the
given class is killed by some power of $\mathcal P$.
Thus $\mathcal P$ is locally nilpotent on the cohomology of every
fixed $Q_D$.

Now form the quotient complex
\[
 Q^\bullet=
 C^\bullet_{\hthreects}(G,B_{\hthreehol})\big/
       \colim_d C^\bullet_{\hthreects}(G,A_d)
   =\colim_{D>C}C^\bullet_{\hthreects}(G,Q_D).
\]
The equality follows from the preceding degreewise exact sequences.
Filtered colimits of vector spaces are exact, so the fixed-quotient
argument makes $\mathcal P$ locally nilpotent on $H^a(Q^\bullet)$.
These groups are finite, using the long exact sequence for this
quotient, \Cref{h3:prop:holomorphic-comparison}, and the separately
proved \Cref{h3:prop:bounded-pole-finiteness}.
Both Frobenius and $\mathcal P$ preserve the bounded-pole cochain
union.  Their descended maps on $Q^\bullet$ satisfy
$\mathcal P\mathrm{Fr}=1$.  Thus $\mathcal P$ is surjective,
hence bijective, on its finite cohomology groups.  A bijective locally
nilpotent endomorphism can act only on the zero group.  This proves
that the first arrow of \eqref{h3:ord:natural-comparison} is a
quasi-isomorphism.  The second is
\Cref{h3:prop:holomorphic-comparison}.

Every $A_d$, $B_{\hthreehol}$, and $B_{\hthreepk}$ has finite point-value
cohomology and a Polish cochain complex.  The family-lifting result
\Cref{h3:co:parameters} identifies the cohomology with a profinite
parameter space $T$ at each such stage with continuous maps from $T$ to
its finite discrete cohomology.  In the filtered union, a continuous
map from compact $T$ to the discrete colimit has finite image, and those
finitely many classes occur at one common stage.  Thus this identification
commutes with the stagewise filtered colimit and gives the comparison
for every profinite $T$.  This argument does not assert that $B_0$, or
its quotient cochain complex, is Polish.
Finally the arrows being proved to be quasi-isomorphisms are the
coefficient inclusions, so they retain products and constants.
Write $\phi_p(f)=f^p$ for this untwisted coefficient Frobenius.
It is distinct from the fixed $q$-power operator $\mathrm{Fr}$ used
with the line retractions above.  The inclusions commute with $\phi_p$,
which is an automorphism on perfected cochains; hence $\phi_p$ is an
isomorphism on each of the compared cohomology groups.
\Cref{h3:prop:ordinary-constants} identifies their endpoint with the
displayed group cohomology.  All constructions are natural under the
semilinear coefficient-field action, so the comparison also descends
to the original constants.
\end{proof}

\begin{corollary}[The ordinary exterior classes]\label{h3:ord:exterior-classes}
For $p\geq5$ the last group in \Cref{h3:thm:ordinary-comparison} is
\[
 \Lambda_k(h,e,d),\qquad |h|=|e|=1,\quad |d|=3.
\]
Here $h$ is the logarithmic class of $P_1=\Z_p^\times$, $e$ is
the determinant logarithm of $\GL_2(\Z_p)$, and $d$ restricts to
the orientation class of $\hthreeSL_2(\Z_p)$.  All three classes have
integral lifts in continuous $\Z_p$-cohomology of these frame groups.
\end{corollary}

\begin{proof}
The group $\hthreeSL_2(\Z_p)$ has dimension three, no element of order
$p$, and trivial adjoint orientation.  Its first mod-$p$ cohomology
vanishes.  To see this directly, let $f:\hthreeSL_2(\Z_p)\to k$ be a
continuous homomorphism.  Choose $a\in\Z_p^\times$ with
$a^2-1$ a unit, possible for $p\geq5$.  Conjugation by
$\operatorname{diag}(a,a^{-1})$ gives
\[
 f\left(\begin{smallmatrix}1&a^2t\\0&1\end{smallmatrix}\right)
   =f\left(\begin{smallmatrix}1&t\\0&1\end{smallmatrix}\right).
\]
Additivity on this root subgroup and invertibility of $a^2-1$
force its restriction to vanish.  The lower root subgroup is treated
by the same displayed conjugation, with $a^{-2}$ in place of $a^2$.
These subgroups generate $\hthreeSL_2(\Z_p)$ by row reduction over the
local ring $\Z_p$, so $f=0$.
Compact group duality now gives $H^2=0$ and $H^3=k$, with the
orientation generator, while $H^0=k$.

The determinant quotient $\Z_p^\times$ has cohomology
$\Lambda_k(e)$: average its tame subgroup and use the two-term
resolution of its procyclic quotient.  Its action on $H^3$ of
$\hthreeSL_2$ is trivial, since conjugation on $\mathfrak{sl}_2$ has
determinant one.  The Hochschild--Serre sequence for
\[
 1\longrightarrow\hthreeSL_2(\Z_p)\longrightarrow\GL_2(\Z_p)
 \xrightarrow{\det}\Z_p^\times\longrightarrow1
\]
has only rows $0,3$ and columns $0,1$.  It has no possible
differential, and its edge map onto the orientation row is surjective.
Choose a lift $d$ of that row.  The product $ed$ is nonzero on
the associated graded, and odd squares vanish because $p$ is odd.
This proves the exterior algebra on $e,d$.  The finite-resolution
K\"unneth calculation with $P_1$ adds $h$.

Integrally the same finite resolutions have finitely generated
$\Z_p$-cohomology.  Integral duality gives the invariant
$\Z_p$-orientation line in degree three for $\hthreeSL_2$.
The preceding elementary-matrix argument also gives
$H^1(\hthreeSL_2,\Z_p)=0$.  Reduction embeds
$H^2(\hthreeSL_2,\Z_p)/p$ into $H^2(\hthreeSL_2,\F_p)=0$;
finite generation and Nakayama's lemma therefore give
$H^2(\hthreeSL_2,\Z_p)=0$.
For clarity, vanishing in the integral degree needed for this reduction
also follows from the finite coefficient models.  The length-three
reduced resolution and the finite filtration by powers of $p$ make
$H^a(\hthreeSL_2,\Z/p^\ell)$ finite and zero for $a>3$ at every
$\ell$.  Cochain reductions are surjective by continuous set-theoretic
sections, and the finite cohomology towers are Mittag--Leffler.  The
inverse-limit cochain sequence therefore gives
$H^4(\hthreeSL_2,\Z_p)=0$.  The coefficient reduction long exact
sequence now proves that the integral orientation surjects onto its
mod-$p$ class.  The integral determinant Hochschild--Serre sequence uses the same
two-term completed resolution for the procyclic quotient, after exact
tame averaging.  It again has no possible differential out of the
orientation row, so that orientation lifts to $\GL_2$.  Normalized continuous
logarithms give integral $e$ and $h$.  Reduction of these choices
gives the three specified mod-$p$ classes.
\end{proof}
